\chapter{Marco Teórico}

El análisis automático de código fuente constituye un área interdisciplinaria que combina conceptos de ingeniería de software, procesamiento de lenguaje natural y aprendizaje automático. En entornos educativos, particularmente en cursos introductorios de programación, el análisis de código se utiliza para evaluar soluciones desarrolladas por estudiantes y detectar posibles similitudes entre programas. Este tipo de análisis está estrechamente relacionado con el problema de la detección de clones de código, el cual ha sido ampliamente estudiado en la literatura de ingeniería de software.

Roy y Cordy describen los clones de código como fragmentos de programas que presentan similitudes significativas en su estructura o funcionalidad, y destacan que su detección es una tarea fundamental para el mantenimiento y análisis de software \cite{Roy2007Survey}. No obstante, este proceso enfrenta múltiples desafíos técnicos debido a la ambigüedad estructural del código y a la existencia de múltiples implementaciones para un mismo algoritmo. Esto hace más complejo distinguir entre una similitud semántica real y coincidencias sintácticas por el uso de convenciones de programación estándar y buenas prácticas, dificultando la identificación de clones funcionales con una alta precisión.

La investigación sobre similitud de código identifica niveles en los que dos programas pueden considerarse similares, siendo la clasificación de Bellon \cite{Bellon2007Comparison} el marco predominante. Los clones de Tipo I representan copias exactas con variaciones mínimas en formato, mientras que los de Tipo II permiten el renombrado de identificadores manteniendo una estructura sintáctica idéntica. Esta taxonomía establece una base estandarizada que permite definir matemáticamente los primeros dos tipos como relaciones de equivalencia precisas \cite{Bellon2007Comparison}. Debido a su naturaleza textual, su detección automática resulta altamente confiable mediante técnicas de normalización de tokens o análisis sintáctico básico.

Por el contrario, los clones de Tipo III y IV introducen una complejidad significativamente mayor para los sistemas de análisis automático. Según detalla Bellon \cite{Bellon2007Comparison}, los de Tipo III aceptan alteraciones donde se añaden o eliminan sentencias de código, lo que vuelve su definición inherentemente subjetiva y vaga. En estos casos, la similitud pierde su propiedad transitiva, dificultando la labor de las herramientas que carecen de la abstracción necesaria para identificar roles sintácticos equivalentes. Estas limitaciones evidencian la brecha existente entre el análisis superficial de secuencias y la comprensión de la lógica subyacente.

Finalmente, los clones de Tipo IV, o clones semánticos, constituyen el desafío analítico más crítico dentro del marco conceptual referido \cite{Bellon2007Comparison}. Estos fragmentos implementan la misma funcionalidad mediante estructuras sintácticas completamente diferentes, compartiendo únicamente precondiciones y postcondiciones similares. La detección automatizada de esta categoría se considera un problema técnicamente complejo e incluso indecidible, ya que requiere interpretar la intención del código en lugar de su forma. Por ello, el análisis del Tipo IV \cite{Bellon2007Comparison} exige métodos que trasciendan la comparación estructural para evaluar la lógica computacional pura.

Uno de los enfoques más utilizados para analizar la estructura de los programas consiste en representar el código mediante árboles de sintaxis abstracta (AST). Un AST es una representación jerárquica que describe la estructura gramatical de un programa de acuerdo con las reglas del lenguaje de programación \cite{Pradeep2025AntiPlagiarism}. Esta estructura se genera mediante un proceso de análisis sintáctico (parsing) empleando herramientas como tree-sitter o ANTLR, las cuales transforman el texto en constructos gramaticales \cite{Song2024AST}. De este modo, el AST retiene la jerarquía de los elementos funcionales como declaraciones y ciclos, mientras ignora detalles léxicos irrelevantes como el formato, los comentarios o los nombres de los identificadores \cite{Pradeep2025AntiPlagiarism}.

Este tipo de representación permite comparar programas a nivel estructural, facilitando la detección de similitudes incluso cuando existen diferencias superficiales. Al cuantificar la coincidencia entre nodos o emplear métricas como la Distancia de Edición del Árbol (TSED), los sistemas pueden identificar arquitecturas lógicas con una precisión cercana al 90\% en casos de plagio o clones con modificaciones estructurales \cite{Pradeep2025AntiPlagiarism}. Investigaciones recientes han demostrado que la distancia entre árboles sintácticos constituye una métrica robusta para estimar la similitud, especialmente cuando el código ha sufrido reordenación de sentencias o tácticas de ofuscación comunes que no alteran la jerarquía gramatical del programa \cite{Song2024AST}.

No obstante, el uso exclusivo de AST presenta limitaciones críticas para identificar similitudes semánticas profundas, como las presentes en los clones de Tipo IV. Existe una divergencia notable entre la forma gramatical y la función, donde dos programas pueden producir árboles casi idénticos, pero ejecutar lógicas funcionales opuestas \cite{Song2024AST}. Además, la comparación de subárboles implica un alto costo computacional y muestra vulnerabilidades ante técnicas de ofuscación extrema o código generado por inteligencia artificial, donde la lógica se reestructura de forma imprevista sin perder su propósito original, superando la capacidad de análisis de un enfoque puramente sintáctico.

Además del análisis sintáctico, algunos enfoques utilizan representaciones basadas en grafos que permiten capturar relaciones más complejas dentro del programa. Entre estas representaciones se encuentran los grafos de flujo de control (CFG) y los grafos de flujo de datos (DFG) \cite{Yu2023Graph}. Un CFG describe las posibles rutas de ejecución, mientras que un DFG representa las dependencias entre variables y operaciones. A diferencia de los árboles de sintaxis abstracta (AST), que se limitan a la jerarquía gramatical, estas representaciones modelan el comportamiento dinámico al explicitar el orden de ejecución y la propagación de valores. Mientras que un AST visualiza las sentencias como bloques secuenciales, un CFG revela la estructura causal subyacente mediante aristas que conectan predicados lógicos con sus operaciones objetivo \cite{Yu2023Graph}.

Estas representaciones permiten modelar el comportamiento del programa y capturar información semántica relacionada con la forma en que los datos se transforman a lo largo de la ejecución. Frecuentemente, el flujo de control y el de datos se integran en un Grafo de Dependencia del Programa (PDG), el cual rastrea la trayectoria de una variable desde su definición hasta su uso final \cite{Lu2026CSSG}. Esta sensibilidad al flujo permite detectar divergencias semánticas drásticas en códigos que, estructuralmente, poseen ASTs idénticos. Al mapear dependencias de datos mediante aristas direccionales, el PDG distingue variaciones funcionales basadas en el origen de los valores, superando la jerarquía plana propia de los enfoques puramente sintácticos \cite{Lu2026CSSG}.

Como señala la literatura, el uso de grafos permite aprender representaciones semánticas del código más ricas que aquellas obtenidas mediante análisis puramente textual. Esto se debe a que los grafos de flujo proporcionan una abstracción global que gobierna la ejecución condicional y el movimiento de valores, superando la granularidad local de los nodos de un árbol sintáctico. Al integrar información sobre el contexto condicional y la causalidad lógica, estas representaciones facilitan la identificación de clones semánticos que, aunque difieran totalmente en su forma gramatical, mantienen una identidad funcional observable en su topología de flujo \cite{Yu2023Graph, Lu2026CSSG}.

Con el avance del aprendizaje automático, se han desarrollado métodos que permiten aprender representaciones del código a partir de grandes colecciones de programas. Uno de los conceptos fundamentales en este contexto es el de representaciones distribuidas, introducido por Mikolov en el desarrollo de modelos de embeddings para lenguaje natural \cite{Mikolov2013Distributed}. En esta teoría, los elementos de una secuencia se representan como vectores dentro de un espacio de alta dimensión donde la proximidad entre vectores refleja similitudes semánticas. En este esquema, el significado de un objeto no se localiza en un único símbolo, sino que se distribuye a través de múltiples componentes numéricos, permitiendo codificar conceptos complejos de manera compacta en vectores continuos \cite{Mikolov2013Distributed}.

Este principio puede aplicarse también al código fuente, permitiendo representar fragmentos de programas como vectores que capturan características relevantes de su estructura y comportamiento. Para una adaptación efectiva, el código no se procesa como texto plano, sino que se descompone en colecciones de caminos sintácticos extraídos del AST. Mediante el uso de mecanismos de atención neuronal, el modelo pondera la importancia de cada trayectoria para consolidarlas en un único vector global o \textit{code embedding} que condensa la funcionalidad total del fragmento \cite{Alon2018Code2Vec}. Este proceso permite que la ubicación de los vectores en el espacio vectorial refleje afinidades funcionales y roles prácticos del código analizado.

La importancia de estas representaciones distribuidas radica en su capacidad para transformar estructuras discretas y complejas en vectores numéricos comparables mediante técnicas matemáticas de similitud, como la similitud de coseno. Este enfoque permite realizar aritmética vectorial para identificar analogías de programación y reduce la complejidad espacial de un nivel polinomial a uno lineal, facilitando el procesamiento de millones de funciones \cite{Alon2018Code2Vec}. Al habilitar esta viabilidad algorítmica, se potencia la capacidad de generalización del modelo frente a estructuras no observadas durante el entrenamiento, superando las limitaciones de rigidez propias de los métodos simbólicos tradicionales \cite{Mikolov2013Distributed, Alon2018Code2Vec}.

Una de las métricas más utilizadas es la similitud del coseno, ampliamente empleada en recuperación de información y análisis de texto. De acuerdo con Manning, Raghavan y Schütze, esta técnica mide el ángulo entre dos vectores dentro de un espacio vectorial, proporcionando una medida de similitud independiente de su magnitud. Matemáticamente, este valor se obtiene mediante el producto escalar de los vectores dividido por el producto de sus longitudes euclidianas, lo que normaliza los elementos a vectores unitarios \cite{Manning2008Introduction}. Este enfoque permite que la métrica se centre exclusivamente en la orientación de las características, evitando que diferencias en la longitud del código o en la frecuencia de términos distorsionen la percepción de afinidad semántica.

En el contexto del análisis de código, esta métrica permite comparar embeddings generados por modelos de aprendizaje automático y determinar qué tan similares son dos fragmentos. Dado que las representaciones distribuidas sitúan conceptos similares en regiones espaciales cercanas, la proximidad angular refleja directamente la identidad funcional del programa \cite{Alon2018Code2Vec}. Esta propiedad es fundamental en entornos de aprendizaje profundo, pues permite que el sistema realice razonamientos analógicos mediante aritmética vectorial. Así, la similitud del coseno facilita la identificación de relaciones semánticas complejas que no fueron explícitamente etiquetadas durante el entrenamiento del modelo.

Finalmente, al aplicar esta métrica sobre modelos como \textit{code2vec}, se hace posible descubrir el propósito funcional del software a través de sus vecinos más cercanos. Por ejemplo, la similitud del coseno permite que el sistema identifique nombres semánticamente equivalentes para funciones, asociando vectores de términos como \textit{size}, \textit{length} o \textit{count} debido a su coincidencia en el espacio vectorial \cite{Alon2018Code2Vec}. Asimismo, permite capturar analogías de programación profundas, donde la combinación vectorial de operaciones distintas, como \textit{equals} y \textit{toLowerCase}, arroja una alta similitud con el vector de \textit{equalsIgnoreCase}, demostrando una comprensión robusta de la semántica del código.

El desarrollo reciente de modelos basados en la arquitectura Transformer ha permitido avances significativos en el análisis automático de texto y código. Los Transformers se fundamentan en el mecanismo de autoatención \cite{Vaswani2017Attention}, el cual permite que cada elemento de una secuencia pueda atender a otros dentro de la misma para capturar relaciones contextuales. Este proceso funciona mediante el mapeo de la entrada hacia vectores de consultas (\textit{queries}), claves (\textit{keys}) y valores (\textit{values}), donde la relevancia de cada componente se determina mediante el producto punto entre consultas y claves. Además, el uso de atención multicabezal permite al modelo procesar información de distintos subespacios de representación en paralelo, logrando una comprensión global de la secuencia \cite{Vaswani2017Attention}.

Este mecanismo permite analizar dependencias de largo alcance dentro de los datos, lo que resulta útil para comprender estructuras complejas como las presentes en el código fuente. A diferencia de las redes neuronales recurrentes (RNN), que requieren un camino de longitud $O(n)$ para conectar elementos distantes, los Transformers reducen esta trayectoria a un número constante de operaciones $O(1)$. Al prescindir del procesamiento paso a paso, la arquitectura establece saltos directos entre cualquier par de posiciones, facilitando el aprendizaje de dependencias remotas que en modelos tradicionales se perderían debido a la longitud de la secuencia \cite{Vaswani2017Attention}.

Dicha capacidad es idónea para el software, donde elementos vinculados lógicamente, como la definición de una función y su correspondiente retorno, suelen estar separados por grandes distancias textuales. La investigación indica que las capas profundas de los Transformers reconstruyen implícitamente el flujo de datos y de control, aprendiendo representaciones jerárquicas del código \cite{Zhang2022TransformerCode}. Asimismo, la especialización semántica de las cabezas de atención permite que ciertos componentes del modelo se dediquen exclusivamente a rastrear el uso de variables o la estructura algorítmica, superando la rigidez de los modelos secuenciales previos \cite{Zhang2022TransformerCode}.

A partir de la arquitectura Transformer han surgido modelos especializados en el análisis de código, siendo CodeBERT uno de los más destacados. Desarrollado por Microsoft, este modelo se basa en una estructura bimodal de múltiples capas bidireccionales similar a RoBERTa-base \cite{Feng2020CodeBERT}. Su funcionamiento se fundamenta en un aprendizaje auto-supervisado que combina lenguaje natural (NL) y de programación (PL) mediante dos tareas críticas: el Modelado de Lenguaje Enmascarado (MLM) y la Detección de Tokens Reemplazados (RTD). Mientras que el MLM entrena al modelo para predecir elementos ocultos basándose en el contexto, la tarea RTD le permite actuar como un discriminador que identifica si un token es original o una alternativa generada, logrando así una comprensión profunda de las distribuciones léxicas en ambos dominios \cite{Feng2020CodeBERT}.

Este modelo utiliza el aprendizaje auto-supervisado para aprender representaciones contextuales del código, lo que le permite capturar relaciones semánticas entre fragmentos de programas. Dicha capacidad se adquiere a través del entrenamiento con millones de funciones de GitHub en seis lenguajes de programación, empleando tanto datos bimodales donde el código se empareja con su documentación, como datos unimodales masivos. Esta estrategia facilita una alineación transversal entre texto y código, permitiendo que el modelo aprenda implícitamente qué conceptos del lenguaje natural corresponden a estructuras sintácticas específicas y facilitando la inferencia de lógica computacional a partir de descripciones funcionales \cite{Feng2020CodeBERT}.

En consecuencia, estas representaciones pueden aplicarse con éxito en tareas como búsqueda de código o análisis de similitud. La utilidad de CodeBERT radica en su habilidad para reconocer la equivalencia semántica de arquitecturas lógicas complejas, identificando funciones que realizan la misma tarea, pero con sintaxis o llamadas a la API divergentes. Mediante un proceso de ajuste fino (\textit{fine-tuning}), este modelo alcanza niveles de recuperación del 96\% en la detección de clones semánticos de Tipo IV \cite{ReplicationCodeBERT2022}. Esto lo consolida como una herramienta de vanguardia para la búsqueda semántica y la identificación de similitudes profundas, superando las capacidades de los métodos basados únicamente en estructuras sintácticas o textuales.

Posteriormente se desarrolló GraphCodeBERT, una extensión de CodeBERT que incorpora información estructural adicional mediante grafos de flujo de datos. Según Guo \cite{Guo2021GraphCodeBERT}, este modelo integra relaciones de dependencia entre variables dentro del proceso de aprendizaje, lo que permite capturar mejor la semántica del programa. Este avance se logra al abandonar la concepción del código como una secuencia plana de tokens para integrar la lógica mediante mecanismos de atención enmascarada guiada por grafos. Específicamente, el modelo emplea una representación de entrada bimodal que concatena comentarios, tokens de código y variables, utilizando tareas de pre-entrenamiento como la predicción de bordes y la alineación de nodos para fortalecer el vínculo entre el flujo de datos y la sintaxis \cite{Guo2021GraphCodeBERT}.

En lugar de analizar únicamente secuencias de tokens, GraphCodeBERT utiliza la estructura del flujo de datos para comprender cómo interactúan las variables dentro del código. El seguimiento de estas dependencias permite rastrear la procedencia de los valores, capturando la intención funcional incluso ante convenciones de nomenclatura deficientes o nombres de variables poco descriptivos. Al modelar la relación causal, el sistema logra desambiguar el rol de identificadores idénticos con funciones distintas y manejar dependencias de largo alcance que los modelos secuenciales suelen omitir \cite{Guo2021GraphCodeBERT}. Esta capacidad garantiza que el análisis se centre en la transformación de los datos, proporcionando una abstracción semántica superior.

Esta arquitectura mejora significativamente la capacidad para identificar programas que implementan la misma lógica mediante estructuras sintácticas diferentes. La idoneidad de GraphCodeBERT para este caso de estudio radica en su enfoque en el comportamiento real del código, ignorando la jerarquía gramatical innecesaria de los árboles sintácticos para priorizar la semántica crítica. Al omitir detalles superficiales e irrelevantes, el modelo identifica con precisión clones semánticos que presentan estilos de programación divergentes \cite{Guo2021GraphCodeBERT, MartinezGil2025Interpretability}. Por tanto, el uso de grafos de flujo de datos permite detectar la identidad lógica en fragmentos donde la ofuscación o las diferencias estructurales invalidarían los métodos basados únicamente en secuencias \cite{MartinezGil2025Interpretability}.

El análisis estilométrico es una disciplina que estudia las características estadísticas del lenguaje utilizadas por un autor. Frantzeskou demostró que es posible identificar autores de código fuente mediante el análisis de patrones de escritura, como el uso de nombres de variables y convenciones de formato \cite{Frantzeskou2006Source}. Dicha identificación se sustenta en la extracción de rasgos categorizados en niveles léxicos, sintácticos, estructurales y semánticos. Mientras que los rasgos léxicos capturan elementos superficiales mediante n-gramas de tokens, las características sintácticas y estructurales analizan la jerarquía de las sentencias y la organización de alto nivel, permitiendo obtener una representación multidimensional del estilo del programador \cite{Frantzeskou2006Source, Horvath2026ProgrammerAttribution}.

Estas características pueden utilizarse como una especie de “huella digital” asociada al estilo de programación de un autor. La efectividad de este enfoque radica en que, a pesar de las restricciones formales de los lenguajes de programación, existe una flexibilidad inherente que permite a los desarrolladores adquirir hábitos consistentes. Estos patrones son el resultado de preferencias personales y cognitivas sobre la estructuración de la lógica y la resolución de problemas, lo que genera comportamientos habituales repetitivos. Diversos estudios han mostrado que los programadores tienden a mantener estos rasgos estables, consolidando una firma estilística que persiste en diferentes contextos de desarrollo \cite{Horvath2026ProgrammerAttribution}.

Para materializar la identificación, el código fuente se transforma en vectores numéricos que facilitan tareas de atribución y verificación de autoría. Mientras que la atribución busca asignar un código a un autor dentro de un conjunto de candidatos, la verificación valida de forma binaria si el fragmento corresponde a la autoría reclamada. Actualmente, este proceso emplea modelos que varían desde algoritmos clásicos, como Random Forests y k-NN, hasta arquitecturas avanzadas de aprendizaje profundo, como redes neuronales convolucionales y Transformers, que logran distinguir patrones estilísticos complejos y jerárquicos directamente de la fuente \cite{Horvath2026ProgrammerAttribution}.

Para analizar automáticamente estos patrones estilísticos, diversos estudios han utilizado redes neuronales convolucionales a nivel de caracteres (CharCNN). En este enfoque, el código se procesa como una secuencia de caracteres individuales en lugar de tokens o palabras. Zhang, Zhao y LeCun \cite{Zhang2015CharCNN} demostraron que las redes convolucionales a nivel de caracteres pueden aprender patrones estructurales sin necesidad de realizar un análisis lingüístico previo. Este proceso se inicia cuantificando el texto letra por letra mediante una codificación ``$1$-de-$m$'', donde cada carácter se representa como un vector dentro de un alfabeto de tamaño $m$ predefinido. Al tratar el código como una señal en bruto con una longitud máxima fija $l_0$, el modelo logra una representación granular que incluye elementos frecuentemente ignorados por los analizadores de tokens, como signos de puntuación, operadores y caracteres de salto de línea \cite{Zhang2015CharCNN}.

Este tipo de redes utiliza filtros convolucionales que detectan patrones locales en la secuencia de caracteres, lo que permite identificar características relacionadas con el estilo del código analizado. A través de convoluciones temporales $1-D$, el modelo emplea núcleos o kernels que se deslizan secuencialmente para identificar combinaciones locales de caracteres. Posteriormente, se aplican módulos de max-pooling temporal para extraer las señales más significativas, permitiendo que la red construya abstracciones profundas a partir de fragmentos de información mínima. Esta arquitectura facilita la detección de rasgos estilísticos sutiles, como preferencias en la disposición de operadores o hábitos de espaciado, que conforman la identidad técnica de un programador \cite{Zhang2015CharCNN}.

La principal ventaja de las CharCNN radica en su capacidad para capturar el estilo sin depender de diccionarios o reglas gramaticales explícitas. Al no requerir segmentación formal, el modelo aprende naturalmente combinaciones atípicas de caracteres y hábitos de codificación idiosincrásicos que se manifiestan en la escritura en bruto. Esta independencia lingüística resulta fundamental para la estilometría de código, ya que permite identificar la ``huella digital'' del autor basándose en micro-patrones estructurales que son intrínsecos a su comportamiento habitual, independientemente de si el lenguaje analizado permite una segmentación formal en palabras o no \cite{CodeStylist2023}.

A partir de estas observaciones, diversos estudios recientes han propuesto arquitecturas híbridas que integran múltiples fuentes de información para obtener representaciones más completas del código fuente. En este enfoque, las características semánticas aprendidas por modelos basados en Transformers, como GraphCodeBERT, se combinan con rasgos estilométricos extraídos mediante redes convolucionales a nivel de caracteres u otros métodos de análisis estadístico. Este proceso, conocido como fusión de características (\textit{feature fusion}) \cite{MartinezGil2026Improving}, permite concatenar o integrar diferentes vectores de representación dentro de un mismo espacio de aprendizaje, de modo que el modelo pueda considerar simultáneamente la lógica funcional del programa y los patrones estilísticos asociados a su escritura. Al aprovechar estas perspectivas complementarias, los sistemas híbridos logran construir representaciones más ricas del código fuente, lo que mejora la capacidad de generalización del modelo y aumenta su precisión en tareas complejas como la detección de similitud semántica, la identificación de clones de código y el análisis de autoría \cite{MartinezGil2026Improving}.

En conjunto, la literatura muestra que el análisis automático de código fuente requiere considerar múltiples niveles de representación, desde estructuras sintácticas y relaciones de flujo hasta patrones semánticos y estilísticos. Los enfoques tradicionales basados únicamente en comparaciones textuales o estructurales resultan insuficientes para capturar la complejidad inherente de los programas modernos. En contraste, los métodos basados en aprendizaje profundo permiten construir representaciones distribuidas capaces de modelar tanto la lógica computacional como los patrones característicos de escritura del código. En este contexto, las arquitecturas híbridas que combinan modelos semánticos basados en Transformers con técnicas de estilometría representan una estrategia prometedora para abordar el problema de la similitud de código desde una perspectiva más integral.

En conjunto, estos fundamentos teóricos permiten construir sistemas de análisis automático de código que integran múltiples enfoques. Por un lado, los modelos basados en Transformers permiten capturar la semántica del código mediante representaciones vectoriales que reflejan la lógica del programa. Por otro lado, las técnicas de análisis estilométrico permiten identificar patrones asociados a la forma en que el código es escrito. La combinación de estos enfoques permite desarrollar sistemas más robustos para analizar programas, ya que integran información tanto sobre la lógica del código como sobre su estilo de escritura. Este tipo de aproximación resulta especialmente relevante en contextos educativos donde se busca evaluar la equivalencia lógica entre programas y analizar posibles indicios de generación automática de código.